\section{SUNIQ and ZUNIQ packaging}

Although they are not used in this standard, the SUNIQ and ZUNIQ packing schemes
are implemented in some MOC libraries and are therefore briefly described here.
They are valid with all quantities.

\subsection{SUNIQ packaging}

SUNIQ provides a generalized alternative to the UNIQ packing, which is specific to HEALPix indices.
Its main drawback is that it requires one additional bit.
SUNIQ relies on a sentinel bit (hence the “S” in SUNIQ).
Given a cell index $index$, the SUNIQ value is obtained by setting a sentinel bit immediately to the
left of the most significant bit (MSB) of the index.
 The resulting expressions for HEALPix and TMOC/FMOC indices are:
\begin{eqnarray*}
  suniq_{hpx} & = & 4^{order + 2} + index \\
  suniq_{t/f} & = & 2^{order + 1} + index
\end{eqnarray*}

The inverse operation uses the {\textit leading\_zeros} function,
which returns the number of leading zeros in the binary representation of a 64-bit (unsigned) integer.
\begin{eqnarray*}
  order_{hpx} & = & \frac{59 - leading\_zeros(suniq)}{2} \\
  index_{hpx} & = & suniq - 4^{order + 2} \\
  order_{t/f} & = & 62 - leading\_zeros(suniq) \\
  index_{t/f} & = & suniq - 2^{order + 1}
\end{eqnarray*}

Sorting UNIQ or SUNIQ values produces the same ordering of cells: cells of lower order appear first,
and within a given order, cells are ordered by increasing index.


\subsection{ZUNIQ packaging}

ZUNIQ encodes both the cell index and its order in a single integer while preserving the
ordering of cells according to their ranges at the maximum resolution.
HEALPix indices correspond to the concatenation of 12 Z-order curves and therefore follow the
natural ordering of a multi-resoltuion Morton (Z-order) curve, which motivates the name ZUNIQ.
ZUNIQ is constructed from the starting index of the cell at the maximum resolution, left-shifted by one bit,
with the bit at the right of the least significant bit used as a sentinel.
This scheme is similar to the numbering used by
\textit{S2CellId} in the Google \textit{S2Geometry} library.
\begin{eqnarray*}
  zuniq_{hpx} & = & 2 \times index_{29} + 4^{29-order}\\
  zuniq_{t}   & = & 2 \times index_{62} + 2^{62-order}\\
  zuniq_{f}   & = & 2 \times index_{52} + 2^{52-order}
\end{eqnarray*}
More generally:
\begin{eqnarray*}
  zuniq_{hpx} & = & 2 \times index_{order_{max}} + 4^{order_{max}-order} \\
  zuniq_{t/f} & = & 2 \times index_{order_{max}} + 2^{order_{max}-order}
\end{eqnarray*}

The inverse transformation uses the {\textit trailing\_zeros} function,
which returns the number of trailing zeros in the binary representation of a 64-bit unsigned integer:
\begin{eqnarray*}
  order_{hpx} & = & 29 - \frac{trailing\_zeros(zuniq)}{2} \\
  index_{hpx} & = & \frac{zuniq - 4^{order_{max}-order}}{2} \\
  order_{t/f} & = & order_{max} - trailing\_zeros(zuniq) \\
  index_{t/f} & = & \frac{zuniq - 2^{order_{max}-order}}{2}
\end{eqnarray*}
Sorting cells according to their ZUNIQ values enables efficient transformations between lists of cells
at mixed resolutions and lists of ranges.
More generally, such orderings allow efficient streamed operations.
